\documentclass[a4paper,12pt]{article}
\usepackage[utf8]{inputenc}
\usepackage[T1]{fontenc}
\usepackage[brazil]{babel}
\usepackage{geometry}
\geometry{margin=2.5cm}
\usepackage{enumitem}
\usepackage{hyperref}
\usepackage{graphicx}
\usepackage{titlesec}

\title{\textbf{Fundamentos do Design Gráfico}\\[4pt]\large Aula 02 --- História e Evolução do Design}
\author{Professor Pedro Sacramento}
\date{16 de junho de 2026}

\begin{document}

\maketitle

% ============================================
\section{Introdução}

Compreender a história do design gráfico é essencial para formar profissionais conscientes das referências visuais que moldaram a comunicação ao longo dos séculos. O design não surge no vácuo: cada estilo, técnica e escola reflete o contexto social, tecnológico e cultural de sua época.


% ============================================
\section{Vídeo}

\noindent Para esta atividade, assista ao vídeo:

\begin{center}
\href{https://www.youtube.com/watch?v=UYhBYu_yauU}{\textit{``The interconnected world of graphic design | Douglass Scott | TEDxRISD''}}
\end{center}

\bigskip
\noindent\textbf{Pergunta 1:} Que áreas o preletor menciona como fontes de inspiração para seu trabalho?

\medskip
\noindent\textit{Resposta:} design gráfico, teatro, dança, pintura, cinema, arquitetura, música, colagem, literatura, ilustração, fotografia.

\medskip
\noindent\textbf{Pergunta 2:} De que forma a perspectiva do preletor pode me ajudar a valorizar meu trabalho no campo do design gráfico?

\medskip
\noindent Nem sempre o valor de uma obra criativa pode ser medido de forma objetiva. \href{https://www.terra.com.br/byte/mais-de-mil-comentarios-criticaram-uma-imagem-de-ia-so-que-era-um-monet-verdadeiro,83718b3ea5d626e0ccf59ee5233b3f8738uffiw1.html}{Um exemplo notável}: em 2026, a pintura original de Claude Monet, \textit{Water Lilies} (1916), foi postada em uma rede social como se tivesse sido gerada por inteligência artificial e recebeu uma enxurrada de críticas --- mais de mil comentários acusando a imagem de ser ``pouco criativa''. A obra era um Monet verdadeiro, quadro comercializado por dezenas de milhões de dólares.

\begin{center}
\includegraphics[width=0.35\textwidth]{monet.jpg}

\medskip
\textit{Water Lilies} (1916), de Claude Monet --- criticada como ``pouco criativa'' ao ser confundida com uma imagem gerada por IA.
\end{center}

\medskip
\noindent O exemplo de Monet mostra como a avaliação das pessoas a respeito do que veem pode ser subjetiva. O valor percebido em uma obra também precisa ser comunicado. Conectar uma criação a diferentes áreas --- teatro, dança, arquitetura, fotografia, literatura, música --- pode ser um bom caminho para ajudar as pessoas a compreenderem a profundidade de algo que criamos, inclusive seu valor financeiro.

\medskip
\noindent De forma prática, você pode começar pequeno e ir expandindo seu alcance:

\medskip
\noindent \textbf{No círculo pessoal:} aniversário de parentes e amigos, peças teatrais em eventos da sua igreja.

\noindent \textbf{Na comunidade local:} aniversário da cidade, feiras de artesanato, festivais gastronômicos locais, eventos esportivos comunitários, celebração de datas comemorativas regionais.

\noindent \textbf{Em instituições e causas:} aniversário de uma universidade, reconhecimento do trabalho de artistas da sua região, homenagem a figuras históricas locais, campanhas de preservação do patrimônio histórico, projetos de revitalização de espaços públicos, enaltecimento de eventos de importância cultural, concursos com viés cultural.

% ============================================
\section{Linha do tempo}

A história do design gráfico pode ser organizada nos seguintes períodos principais:

\begin{itemize}[nosep]
  \item Pré-história e primeiros sistemas de escrita (pictogramas, cuneiforme, hieróglifos)
  \item Manuscritos medievais e a invenção da imprensa (Gutenberg, séc. XV)
  \item Renascimento e expansão da tipografia
  \item Revolução Industrial e produção em massa (séc. XVIII--XIX)
  \item Movimento Arts and Crafts (William Morris, final do séc. XIX)
  \item Art Nouveau e o cartaz publicitário
  \item Bauhaus e o modernismo (1919--1933)
  \item Estilo Tipográfico Internacional / Design Suíço (décadas de 1950--70)
  \item Pós-modernismo e a ruptura das regras (décadas de 1980--90)
  \item Revolução digital e design contemporâneo
\end{itemize}

% ============================================
\section{Exercício}

\subsection*{Objetivo}
Construir coletivamente um panorama da história do design gráfico por meio de pesquisa sobre os períodos históricos da linha do tempo.

\subsection*{Dinâmica}
\begin{enumerate}[nosep]
  \item Os alunos se organizarão em grupos de 1 ou mais pessoas.
  \item Cada grupo ficará responsável por \textbf{um período histórico} da linha do tempo acima.
  \item O grupo fará uma pesquisa sobre o período. Uma sugestão de fonte é o livro aberto do OpenTextBC: \url{https://opentextbc.ca/graphicdesign/}. O grupo produzirá um resumo estruturado contendo:
  \begin{itemize}[nosep]
    \item Contexto histórico (o que estava acontecendo no mundo na época)
    \item Principais características visuais do período
    \item Designers, artistas ou movimentos representativos
    \item Inovações técnicas ou tecnológicas relevantes
    \item Exemplos visuais marcantes
  \end{itemize}
  \item Ao final, cada grupo apresentará seu resumo para a turma, montando juntos a linha do tempo completa.
\end{enumerate}

\bigskip
\noindent O resultado final será um panorama completo da história do design gráfico, construído de forma colaborativa a partir de pesquisa.

% ============================================
\section{Referências}

\begin{itemize}[nosep]
  \item Livro-texto aberto: \href{https://opentextbc.ca/graphicdesign/}{Graphic Design and Print Production Fundamentals} (OpenTextBC, licença Creative Commons)
  \item Vídeo TEDxRISD: \href{https://www.youtube.com/watch?v=UYhBYu_yauU}{\textit{The interconnected world of graphic design | Douglass Scott}}
  \item Reportagem: \href{https://www.terra.com.br/byte/mais-de-mil-comentarios-criticaram-uma-imagem-de-ia-so-que-era-um-monet-verdadeiro,83718b3ea5d626e0ccf59ee5233b3f8738uffiw1.html}{Mais de mil comentários criticaram uma imagem de IA --- só que era um Monet verdadeiro} (Terra, 2026)
\end{itemize}

% Fim da Aula 02

\end{document}
